\documentclass[11pt]{article}
\usepackage[a4paper,margin=2.4cm]{geometry}
\usepackage{amsmath,amssymb}
\usepackage{booktabs}
\usepackage{tikz}
\usetikzlibrary{arrows.meta,positioning,calc,fit,backgrounds}
\usepackage{times}
\usepackage[hidelinks]{hyperref}
\usepackage{caption}
\captionsetup{font=small,labelfont=bf}

\title{\vspace{-1.2cm}Differential Fault Analysis of the Ascon-128 Finalization:\\
How the Attack Works, and Why It Splits the Key in Half}
\author{}
\date{}

\begin{document}
\maketitle
\vspace{-1.4cm}

\begin{abstract}
\noindent
This note explains, step by step, the differential fault attack our tool runs
against the Ascon-128 finalization. A single-bit fault at the input of the last
permutation round leaks the S-box output differential through the tag, because
the key cancels in the difference. The honest result of this fault placement is
sharp and one-sided. Tag-only observation recovers exactly one $64$-bit key
half, $K_1$, with certainty, and it provably cannot touch the other half $K_0$.
Full $128$-bit recovery is possible only when one extra state word is
observable, which the open Vortex SimX substrate exposes but a real
authenticated-encryption device does not. All values below are reproduced by
\texttt{ascon\_key\_recovery.py} and independently re-derived from the Ascon
S-box.
\end{abstract}

\section{Where the key enters: the Ascon-128 finalization}
Ascon keeps a $320$-bit state in five $64$-bit lanes, written
$s = (s_0,s_1,s_2,s_3,s_4)$. The permutation $p^{12}$ applies twelve rounds.
Each round adds a constant, applies a $5$-bit S-box in parallel across the
$64$ bit-columns, then mixes each lane with two rotations. In the finalization
the key is folded in, the final permutation runs, and the tag is read from the
bottom two lanes:
\begin{equation}
T \;=\; \underbrace{\big(p^{12}(S_{\text{in}})\big)_3 \oplus K_0}_{T_0}
   \;\big\|\;
        \underbrace{\big(p^{12}(S_{\text{in}})\big)_4 \oplus K_1}_{T_1}.
\label{eq:tag}
\end{equation}
The key touches the output through a plain XOR, after the last nonlinear layer.
That single fact is what the attack exploits.

\section{The fault model}
The attacker flips one bit at the input $U$ of the \emph{last} round of the
final permutation. The bit sits in a known lane and a known column, so the
injected column difference is $e = 1 \ll (4-\text{lane})$. Write the last round
as
\begin{equation}
S_{\text{out}} \;=\; L\big(\mathrm{Sbox}(\mathrm{AddRC}(U))\big),
\label{eq:lastround}
\end{equation}
where $L$ is the per-lane linear layer. $L$ acts on each lane on its own through
two rotations, so it is a bijection on $\mathrm{GF}(2)^{64}$ and its inverse
$L^{-1}$ is known. Two runs, one clean and one faulted, give the tags $T$ and
$T'$.

\section{Why the key cancels}
Subtract the two tags. In Equation~\eqref{eq:tag} the key is a constant XOR, so
it drops out of the difference:
\begin{equation}
T \oplus T' \;=\; (S_{\text{out}} \oplus S_{\text{out}}')_{\text{lanes }3,4}.
\end{equation}
The attacker now sees the true S-box output differential on lanes $3$ and $4$,
with no knowledge of the key. Applying $L^{-1}$ per lane peels off the linear
mixing and exposes, for every column $j$, two bits of
$\mathrm{Sbox}(x_j) \oplus \mathrm{Sbox}(x_j \oplus e)$, where $x_j$ is the true
$5$-bit S-box input at that column. The S-box difference distribution table then
lists the candidate inputs that could produce those two output bits.

\section{The one-sided outcome: $K_1$ is free, $K_0$ is locked}
Here the Ascon S-box does something specific, and it decides the whole attack.
We verified the following exhaustively over all inputs and differences, and the
independent review re-derived the S-box from first principles before checking
it.

\begin{itemize}
\item The lane-$3$ and lane-$4$ output differential is \emph{invariant} under
      flipping the input on lane $2$, that is under $x \mapsto x \oplus 4$.
      No single-bit mask breaks this, so DDT filtering pins each column only to
      the pair $\{x_j, x_j \oplus 4\}$.
\item The lane-$4$ output bit itself does not depend on the lane-$2$ input bit.
      So $K_1$ is fully determined from the tag alone.
\item The lane-$3$ output bit flips one-for-one with the unknown lane-$2$ input
      bit. So $K_0 = (\text{known}) \oplus L_3(m)$ with $m$ an unknown $64$-bit
      mask. Since $L_3$ is a bijection, $m$ ranges over all $2^{64}$ values, and
      fresh nonces inject a fresh independent $m$ each time.
\end{itemize}
The consequence is a clean split. Tag-only faults at the last round give $K_1$
with certainty and leave an irreducible $2^{64}$ residual on $K_0$. To close the
gap the attacker needs one observable that depends on the lane-$2$ input. The
unmasked capacity word $S_{\text{out},2}$ is exactly such an observable, because
its lane-$2$ output always differs between $x$ and $x \oplus 4$. Reading it
resolves every column and yields the full state, hence $K_0$.

\input{ascon_dfa_figure}

\section{The recovery, end to end}
The pipeline of Figure~\ref{fig:asconflow} is short.
\begin{enumerate}
\item Inject single-bit faults at known columns of the last-round input and
      record each faulty tag.
\item For each column, invert $L$ on the lane-$3,4$ tag difference and read the
      two output-difference bits.
\item Intersect the DDT candidate sets over the masks at that column, which
      leaves the residual pair $\{x_j, x_j \oplus 4\}$.
\item Take the lane-$4$ output bit, which is constant on the pair, and assemble
      $K_1 = L_4(Y_4) \oplus T_1$.
\item For full recovery only, read the capacity word $S_{\text{out},2}$, use it
      to select the correct member of each pair, rebuild the state, and read
      $K_0 = S_{\text{out},3} \oplus T_0$.
\end{enumerate}

\section{Values}
Table~\ref{tab:results} lists the measured cost. The tool validates against the
Ascon known-answer test, then recovers over twenty random keys. The
minimum-fault sweep shows the split directly. Two single-bit masks per column
already give $K_1$ at $100\%$, and no number of single-bit masks gives the full
key without the capacity word.

\begin{table}[t]
\centering
\caption{Measured recovery cost. K1-only is a pure tag-only attack. Full needs
one capacity-word readout in addition to the faults.}
\label{tab:results}
\begin{tabular}{lccc}
\toprule
\textbf{Target} & \textbf{Observable} & \textbf{Faults} & \textbf{Success (20 keys)} \\
\midrule
$K_1$ (64 bits)        & tag only              & $128$ & $100\%$ \\
Full key (128 bits)    & tag $+$ $S_{\text{out},2}$ & $192$ & $100\%$ \\
\bottomrule
\end{tabular}

\vspace{0.5em}
\begin{tabular}{cccc}
\toprule
\textbf{Masks / column} & \textbf{Faults / key} & \textbf{$K_1$ success} & \textbf{Full success} \\
\midrule
$1$ & $64$  & $0\%$   & $0\%$   \\
$2$ & $128$ & $100\%$ & $0\%$   \\
$3$ & $192$ & $100\%$ & $100\%$ \\
\bottomrule
\end{tabular}
\end{table}

\noindent\textbf{Worked example.} With a random key and a fixed nonce,
\begin{quote}
true key \;=\; \texttt{4420823cfde6f1c26b30f90ec7dd01e4}\\
recovered \;=\; \texttt{4420823cfde6f1c26b30f90ec7dd01e4}\quad(match).
\end{quote}
The tag-only path returns $K_1 = \texttt{6b30f90ec7dd01e4}$ and reports $K_0$ as
unrecoverable, exactly as the theory predicts. The full path, given the capacity
word, returns the whole key. Solve time is about $10$\,ms per key.

\section{Scope, stated plainly}
Three limits belong next to every headline number.
First, the honest tag-only result is half the key. Full recovery is contingent
on reading an unmasked capacity word, which the open Vortex SimX substrate can
dump but a fielded AEAD device never outputs. Second, the attack assumes
bit-precise faults at known lane and column at the last-round input. Third, the
$2^{64}$ residual on $K_0$ is specific to last-round, tag-only faults. Faults in
earlier rounds propagate through more nonlinearity and recover more, which is
the setting of prior Ascon fault work. On the real SimX campaign the fault
positions are not yet logged and only the tag is dumped, so the campaign is
shown to produce the right structural class of differentials rather than to
drive an end-to-end recovery.

\end{document}
